\section{Introduction}\label{sec:intro}

In 1914 Ramanujan listed seventeen rapidly converging series for $1/\pi$.
One of them is
\begin{equation}\label{eq:classical}
\sum_{n=0}^{\infty}\frac{(1/2)_n^3}{(n!)^3\,64^n}\,(42n+5)=\frac{16}{\pi},
\end{equation}
where $(a)_0=1$ and $(a)_n=a(a+1)\cdots(a+n-1)$ for $n\ge1$; this is recorded
as \cite[(1.3)]{Au}. A \emph{$q$-analogue} of such a formula is a series in an
extra variable $q$ whose terms tend, after a suitable normalization, to the
terms of the classical series as $q\to1$, and which can itself be summed in
closed form.

Au \cite{Au} found a systematic way to produce $q$-analogues. He starts from a
\emph{Wilf--Zeilberger seed}, a hypergeometric expression from which a whole
family of Wilf--Zeilberger pairs (see \cref{sec:wz}) can be generated. Many of
his $q$-analogues evaluate to infinite products, which are modular forms. Some
evaluate instead to \emph{mock} modular forms, functions that are modular only
after a non-holomorphic correction term is added (see \cref{sec:modular}).
After his Example~XXI, Au writes down a further $q$-analogue of
\eqref{eq:classical}, built from a seed related to Watson's transformation,
and asks whether it too is mock modular \cite[remark after Example~XXI]{Au}.
This paper answers that question. We evaluate the series in closed form,
identify its mock modular part with the one already present in Example~XXI,
give the completion explicitly, and verify the classical limit.

\subsection{The series}\label{sec:series}
For a complex number $a$, a base $r$ with $0<|r|<1$ and an integer $n\ge0$,
put
\[
\poch{a}{r}{n}=\prod_{j=0}^{n-1}(1-ar^j),\qquad
\poch{a}{r}{\infty}=\prod_{j=0}^{\infty}(1-ar^j),
\]
with the empty product equal to $1$. For several parameters we write
$\poch{a_1,\dots,a_s}{r}{n}=\prod_{j}\poch{a_j}{r}{n}$. Throughout, $q$ is a
complex number with $0<|q|<1$, and we use the abbreviations
\[
Q=q^2,\qquad J_d=\poch{q^d}{q^d}{\infty}\quad(d=1,2,3,\dots).
\]
Let $P(q,x)$ be the polynomial
\begin{equation}\label{eq:P}
\begin{aligned}
P(q,x)={}&1+qx+(1+q^2)x^2+2q^3x^3+(3q^2+2q^4)x^4+q^3x^5\\
&+(-2q^2+q^4+q^6)x^6+(-q-q^3+2q^5)x^7-q^4x^8\\
&-(2q^3+3q^5)x^9-2q^4x^{10}-(q^5+q^7)x^{11}-q^6x^{12}-q^7x^{13}.
\end{aligned}
\end{equation}
For each integer $n\ge0$ set
\begin{equation}\label{eq:Bn}
B_n=\frac{\poch{q}{Q}{n}^4\,\poch{Q}{Q^2}{n}^2\,\poch{-Q}{Q}{4n}}
{\poch{Q^2}{Q^2}{2n}^3},\qquad
a_n=\frac{q^{2n}\,P(q,q^{2n})\,B_n}
{(1+q^{8n})(1+q^{2n+1})(1+q^{4n+2})^3},
\end{equation}
and let
\[
S(q)=\sum_{n=0}^{\infty}a_n .
\]
Separating even and odd factors (\cref{fact:evenodd}) shows that $B_n$ can also be
written as
\[
B_n=\frac{\poch{q}{q^2}{n}^4\,\poch{q^2}{q^2}{2n}\,\poch{q^2}{q^4}{n}\,
\poch{q^4}{q^4}{4n}}
{\poch{q^2}{q^2}{4n}\,\poch{q^4}{q^4}{n}\,\poch{q^4}{q^4}{2n}^3},
\]
which is the form in which Au writes the summand. The first term is
\begin{equation}\label{eq:a0}
a_0=\frac{1-q^5}{(1+q)(1+q^2)^2}.
\end{equation}

\subsection{Results}
The main result evaluates $S(q)$.

\begin{theorem}\label{thm:main}
For every complex $q$ with $0<|q|<1$ the series $S(q)$ converges absolutely,
and
\begin{equation}\label{eq:main}
S(q)=\frac{\poch{q}{q^2}{\infty}^2}{J_4^4}\,L(q),\qquad
L(q)=\sum_{r\in\Z}\frac{q^{4r(r+1)}}{1-q^{4r+1}}.
\end{equation}
Both sides extend holomorphically to $q=0$, where they equal $1$.
\end{theorem}

The bilateral series $L(q)$ is the source of mock modularity. Since $L$ also
appears in Au's evaluation of Example~XXI, the two $q$-series are related by an
infinite product (\cref{cor:XXI}).

The next result confirms that $S$ is a $q$-analogue of \eqref{eq:classical};
its proof does not use \cref{thm:main}.

\begin{proposition}\label{prop:classical}
As $q\to1$ through real values $0<q<1$,
\[
\frac{8\,S(q)}{1-q}\longrightarrow\frac{16}{\pi}.
\]
\end{proposition}

Finally, write $q=e^{2\pi i\tau}$ with $\tau$ in the upper half-plane $\HH$.
In \cref{sec:completion} we define an eta quotient $A(\tau)$ of weight $-3/2$,
a function $f(\tau)$ built from Zwegers' $\mu$-function, its completion
$\fh(\tau)=f(\tau)+\frac14R_0(\tau)$ (where $R_0$ is an explicit
non-holomorphic function), and the product $\Sh=A\fh$.

\begin{theorem}\label{thm:completion}
For every $\tau\in\HH$ we have $S(e^{2\pi i\tau})=A(\tau)f(\tau)$. The completed
function $\Sh=A\fh$ satisfies
\[
\Sh\Big(\frac{a\tau+b}{c\tau+d}\Big)=\chi_{-4}(d)\,(c\tau+d)^{-1}\,\Sh(\tau)
\qquad\text{for every }\begin{pmatrix}a&b\\c&d\end{pmatrix}\in\Gamma_0(4),
\]
and the weight-$1/2$ factor $\fh$ is a harmonic Maass form with shadow
\[
\xi_{1/2}\fh=-\tfrac12\,\eta(2\tau)^3 .
\]
In particular $S$ is a mixed mock modular form of weight $-1$.
\end{theorem}

All terms used here ($\Gamma_0(4)$, $\chi_{-4}$, $\xi_k$, harmonic Maass form,
shadow, mixed mock modular form) are defined in \cref{sec:modular}.

\subsection{How the proof works}
The evaluation proceeds in five steps, each short once the right object is
written down.
\begin{enumerate}[label=(\arabic*)]
\item \emph{A finite identity} (\cref{lem:telescope}). We give a two-variable
term $F(n,k)$ and an explicit rational function, the \emph{certificate},
that together form a Wilf--Zeilberger pair. Summing the pair over $k$ shows
that the partial sum $a_0+\dots+a_{N-1}$ equals minus a finite sum $T_N$ of
$2N$ terms.
\item \emph{Watson's transformation} (\cref{lem:watson}). The finite sum $T_N$ is a
terminating very-well-poised series, and Watson's transformation rewrites it as
a sum of $2N$ terms $h_N(k)$ that are all nonnegative when $0<q<1$.
\item \emph{A centred limit} (\cref{lem:centre}). Writing $k=N+r$ and letting
$N\to\infty$, each $h_N(N+r)$ tends to a constant multiple of an explicit term
$b_r$. A bound that does not depend on $N$ justifies exchanging limit and sum.
The result expresses $S(q)$ through the two-sided series $\sum_{r\in\Z}b_r$.
\item \emph{A classical identity} turns $\sum_r b_r$ into the Lambert series
$L(q)$, which proves \cref{thm:main} for $0<q<1$; analytic continuation
gives all $|q|<1$.
\item \emph{Zwegers' theory.} The quotient of $L(q)$ by a theta function is a
special value of Zwegers' $\mu$-function, whose completion and transformation
laws are known. This gives \cref{thm:completion}.
\end{enumerate}
\Cref{prop:classical} is proved separately in \cref{sec:classical} by taking
termwise limits under a uniform bound. \Cref{sec:campbell} compares the
summand with the one in Campbell's $q$-analogue of \eqref{eq:classical}.

\subsection{Computer verification}
Every identity in this paper was checked by computer, using exact rational
arithmetic wherever possible: \cref{thm:main} and \cref{cor:XXI} as
power-series identities modulo $q^{220}$; the certificate identity of
\cref{app:cert} by full polynomial expansion; \cref{lem:telescope,lem:watson,lem:centre}
at many rational values of $q$; and the transformation law and shadow of
\cref{thm:completion} numerically. The scripts are available at
\url{https://github.com/ramanujamath/au-watson-mock-series}.
